Ferroelectricity in doped HfO$_2$ and especially in \hzo{} has reopened the path to scalable ferroelectric memory compatible with CMOS and BEOL processing~\cite{boscke2011,muller2012,schroeder2022}. In a conventional ferroelectric capacitor, however, reading the stored state is destructive because the switched charge $2P_r$ must be sensed by driving the device through its coercive field. A capacitive read avoids that penalty. The two remanent polarization states occupy different small-signal permittivity branches, so a sub-coercive capacitance measurement can distinguish the stored bit without forcing a rewrite~\cite{luo2020}. That possibility is attractive for read-dominated memory operation and for analog or in-memory computing schemes that repeatedly interrogate the same stored state.

Prior work established that a capacitive memory window can be opened and enlarged by interface engineering, and that non-destructive readout can be demonstrated in BEOL-compatible ferroelectric capacitors and array-relevant circuits~\cite{mukherjee2023iedm,lee2025aem}. What those reports do not resolve is the physics that governs the window itself. A large remanent polarization is necessary to define two stored states, but it is not sufficient to guarantee a useful zero-bias CMW: in an ideally symmetric stack, the two $C$--$V$ branches would still intersect at zero bias. The practical questions are therefore sharper than simply asking whether ferroelectricity is present. What determines the usable read bandwidth? Why does the zero-bias window vary between stacks? And will the same readout survive across temperature?

The full TMAT paper addresses those questions by testing two competing pictures. In a ferroelectric-limited picture, the CMW should follow the intrinsic dielectric and switching dynamics of the \hzo{} layer and therefore persist on the ferroelectric timescale. In an interface-limited picture, the window should relax with the impedance of the interfacial region and should change strongly with stack design and temperature. This conference version retains the decisive evidence for the latter view. It develops the argument around one statement: polarization stores the bit, but the interface sets whether and how fast the two states remain capacitively distinguishable.
